\documentclass[10pt]{article}
\usepackage[margin=0.85in]{geometry}
\usepackage{amsmath,amssymb}
\usepackage[hidelinks]{hyperref}

\begin{document}
\begin{center}
  {\Large\bfseries A tree with six-strand braid group of cohomological dimension at least three\par}
  \smallskip
  Disproof of Conjecture 00000004101
\end{center}

Let $\Gamma$ have the path $0-1-2-3-4-5-6$, leaves $7,8$ attached
to $0$, leaf $9$ attached to $3$, and leaves $10,11$ attached to $6$.
We write $P_6(\Gamma)$ and $B_6(\Gamma)$ for the fundamental groups
of its ordered and unordered six-point configuration spaces.

\section*{Acyclicity and the three essential vertices}

Every nonzero vertex has exactly one neighbor of smaller index. Repeatedly
following this neighbor reaches $0$, so the graph is connected. In a
simple cycle, the vertex of largest index would have two distinct lower
neighbors on that cycle. Hence there are no cycles at all and $\Gamma$ is
a tree. Its only vertices of valence at least three are $0,3,6$, each of
valence three. Their closed stars have vertex sets
\[
 \{0,1,7,8\},\qquad \{2,3,4,9\},\qquad \{5,6,10,11\},
\]
which are pairwise disjoint. The open stars are contractible.

\section*{A rank-three subgroup}

Apply Knudsen's toric subgroup construction~\cite[Proposition~3.4(1)]{knudsen}
to $W=\{0,3,6\}$ and six particle labels, assigning the pairs $\{1,2\}$,
$\{3,4\}$, and $\{5,6\}$ to the three vertices, respectively. This is a
binary $W$-partition. For Knudsen's basepoint of bivalent vertices,
subdivide the two edges $1-2$ and $4-5$ once each to obtain six distinct
bivalent vertices. This leaves the three
closed stars unchanged and gives homeomorphic configuration spaces.

For this partition $\lambda$, Proposition~3.4(1) says that the composite
$\delta_\lambda\circ\tau_\lambda$ is injective. Therefore the toric map
$\tau_\lambda\colon\mathbb Z^W\to P_6(\Gamma)$ is injective. The ordered
configuration space covers the unordered one by forgetting the labels,
which induces an injection $P_6(\Gamma)\hookrightarrow B_6(\Gamma)$.
Thus $B_6(\Gamma)$ has a subgroup isomorphic to $\mathbb Z^3$.

Cohomological dimension cannot increase when passing to a subgroup:
restricting a projective resolution over the larger group ring gives a
projective resolution over the subgroup ring, since the larger group
ring is free over the subgroup ring. The three-torus is a classifying
space for $\mathbb Z^3$ and has nonzero third integral cohomology, so
$\operatorname{cd}(\mathbb Z^3)=3$. Consequently
\[
 \operatorname{cd}\bigl(B_6(\Gamma)\bigr)\ge3.
\]
The graph contains no two vertex-disjoint cycles, while its six-strand
braid group violates the asserted upper bound of two.

\section*{Lean verification}

The Lean 4 project checks the finite graph, proves connectivity and
acyclicity from the unique lower neighbors, and checks the essential
vertices, disjoint closed stars, and particle partition. From the
explicitly stated subgroup and dimension assumptions above, it deduces
the contradiction to the proposed dimension criterion. Run
\texttt{lake build} in the \texttt{lean} directory to check it.

\begin{thebibliography}{1}
\bibitem{knudsen}
B.~Knudsen, \emph{The topological complexity of pure graph braid groups is
stably maximal}, \href{https://arxiv.org/abs/2206.06268}{arXiv:2206.06268}.
\end{thebibliography}

\end{document}
